% Reusable results only. Detailed proof: gten18:no-support in arithmetic paper.
\subsection{\texorpdfstring{$G_{5,10}$}{G5,10} の排除から再利用する方法}
\label{ngten18:reusable}

$G_{5,10}=\{01,02,03,12,13,24,34\}$ は，分母の上限なしに排除された．
算術稿の定理 \texttt{gten18:no-support} を参照する．
核心は，三高分母 core の正符号が完全 matching となる定理を，
二価 core の二つの低辺の node 予算と組み合わせることである．
対象は指定した五次の算術的接着模型であり，任意の fibration の分類ではない．

頂点 $4$ の二高分母標準形を $(A,B,C,D)=(2,3,0,1)$ とし，
$a=sS+1,b=zT+1,x=sS/r,y=zT/(r+1)$，$\delta=T-S$ とする．
$q_{23}=1$ なので正の枝であり，$\delta\ge2$ である．
低辺 $04,14$ の node 数をそれぞれ $(n_1,n_2,n_3)$，$(m_0,m_2,m_3)$ とすると，
隣接する三高分母 cores の低 rows が
\[
 R^0=(1-n_1,1-r-n_2,r+2-n_3),\quad
 R^1=(\delta+2-m_0,S+1-m_2,1-T-m_3)
\]
に固定される．各 row には正成分がちょうど一つあり，各 node 組の和は五である．
共通低 puncture $(01;4)$ の三つの接着型は次のように処理する．

\begin{enumerate}
\item 通常型では，有限な node 予算と matching 性だけで
$(r,S,T)=(2,3,5)$ が強制される．整数 pair equation は
$(u-2)(v+1)=-3$ に帰着し，唯一の中央模型は
$(r,s,z,S,T)=(2,2,6,3,5)$，$(a,b)=(7,31)$ である．
しかし反対側の core では，$30$ の signed divisor が
$27\pmod{31}$ であることを要求される．その候補 $27,-4$ はどちらも不適格である．
この部分は $D$ 相殺を使わない．

\item 正例外型では $R^0=(1,-1,-1)$ となる．二つの低辺の $D$ 相殺から
$x-y>0$ と $s-z\le1$ が従う．正の枝の不等式は $s-z\ge1$ を与えるので
$s=z+1$ である．反対側の低 column の二つの負成分がともに $-1$ となり，
$q_{12}+q_{13}=q_{01}+2$ を要求するが，奇分母の parity と矛盾する．

\item 負例外型では $q_{01}\in\{3,5\}$，$\delta=2$ となる．
$h=S-r$，$\alpha=sh/r$，$\beta=z(h+1)/(r+1)$ とすると
\[
 \left(\frac{r^2}{h}-\frac{(r+1)^2}{h+1}\right)\alpha\beta
   =\alpha-\beta+1.
\]
第三角の signed-divisor clutching により，$\alpha,\beta>1$ なら括弧が整数で，
その絶対値は一未満となる．零の場合は $h=r^2/(2r+1)$ となり不可能である．
$\alpha=1$ の例外も保持して処理すると，必ず $\beta=1$ となる．
残る equation は $(h+1)z(z-2)=h-1$ であり，中央模型全体が
\[
 r=3,\quad S=4,\quad T=6,\quad z=2,\quad s=3\alpha,\quad
 a=12\alpha+1,\quad b=13
\]
という無限共鳴族へ帰着する．これを探索上限で切らない．
低辺 $14$ の $D$ と奇分母から $\alpha$ は正の偶数となり，
$q_{12}\ge15\alpha+1$ を得る．低辺 $04$ から $q_{02}\ge8\alpha-7$ を得る．
辺 $24$ の node と matching 性は二つの正整数 $M,N$ に
\[
 M+N=24\alpha+1,\quad M>3q_{02},\quad N>q_{12}
\]
を要求するため，$M+N>39\alpha-20\ge24\alpha+10$ と矛盾する．
\end{enumerate}

再利用すべき点は，低辺の有限な node 予算から隣接三高 core の低 row を固定すること，
matching によってその符号を制限すること，そして第三角の signed divisors で
共鳴族を特定してから $D$ と node の不等式で全族を排除することである．
検算の入口は \nolinkurl{verify_g5_10_obstruction_v18.py} であり，
係数ごとの恒等式二十四本と，証明由来の有限な node の場合分けを保持する．
有界分母探索を一般証明の代わりにしていない．
